% !PRINTED 410
\weberpage{410}
%<<CONT
Functionen constant sein können, so ist nach (6) der Satz 4.
bewiesen.

Diesen Satz wenden wir auf die in 3. vorkommende Function
 $\vartheta$ an. Wenn wir mit den Variablen $x$ irgend eine Permutation
 der alternirenden Gruppe vornehmen, so erfahren die
Functionen $u, u_1, \ldots, u_{\nu-1}$ gleichfalls eine Permutation. Nach
(5) geht $\vartheta$ durch diese Permutation in eine andere Function $\vartheta_1$
über, die nach (4) rational durch $\vartheta$ darstellbar ist. Ebenso ist
aber auch $\vartheta$ rational durch $\vartheta_1$ darstellbar, da man in dem
Satze 3. $\vartheta$ durch $\vartheta_1$ ersetzen kann, und auch $\vartheta$ aus $\vartheta_1$ durch
eine Permutation der alternirenden Gruppe entsteht. Es sind
also nach 4. $\vartheta$ und $\vartheta_1$ \gesperrt{lineare Functionen} von einander.
Daraus ergiebt sich das folgende Resultat:

5.~\gesperrt{Wenn bei einer allgemeinen Gleichung $n^{\text{ten}}$
Grades eine Resolvente mit einem Parameter
besteht, so giebt es eine rationale Function $\vartheta$
von $n$ Variablen $x$, die durch die Permutationen
der alternirenden Gruppe in eine lineare Function
 von sich selbst}
\weq{7}{\chi(\vartheta) = \frac{a\vartheta + b}{c\vartheta + d}}
\gesperrt{übergeht, worin $a, b, c, d$ Constante, d. h. Zahlen
sind.}

\begin{center}
\S.\,109.

\gesperrt{Gruppe der Resolventen mit einem Parameter.}
\end{center}

Die Function $\vartheta$, die im Satze 5. des vorigen Paragraphen
vorkommt, ist nach dem Satze 3. selbst die Wurzel einer Resolvente
 mit \gesperrt{einem Parameter} $\Phi(\vartheta, z) = 0$. Diese Gleichung
ist, da jede Wurzel rational durch jede andere ausdrückbar ist,
eine Normalgleichung, und ihre \gesperrt{Galois}'sche Gruppe ist isomorph
 mit der Gruppe der linearen Substitutionen $\chi(\vartheta)$ des
Satzes 5. Da man nun in der identischen Gleichung $\Phi(\vartheta, z) = 0$
alle Permutationen der alternirenden Gruppe $A$ der Variablen $x$
ausführen kann, wodurch sich $z$ nicht ändert, während $\vartheta$ in jede
andere Wurzel von $\Phi$ übergeht, so ist die Gruppe der linearen
Substitutionen $\chi(\vartheta)$, d. h. die \gesperrt{Galois}'sche Gruppe der Glei-
%CONT>>
